\section{Fazit und Ausblick}
\label{sec:fazit}

\subsection{Zusammenfassung}
SignalReport ist eine vollständige Internet-Qualitäts-Monitoring-Anwendung, die folgende Features umsetzt:

\begin{itemize}
    \item Kontinuierliche Messung von Ping, DNS und HTTP-Latenz
    \item Web-basierte Visualisierung mit Live-Charts und Heatmap
    \item Statistische Auswertung (95. Perzentil, Jitter, Paketverlust)
    \item PDF-Export für Berichte (inkl. 12-Monats-Berichte für Provider)
    \item CSV-Export für Excel-Analysen
    \item IP-Tracking für dynamische IP-Adressen
    \item Konfigurierbare Messziele und Maintenance-Fenster
    \item Challenge-Response-Authentifizierung (SHA-256) für sicheren Zugriff
    \item Setup-Wizard für benutzerfreundliche Ersteinrichtung
    \item Crash-resistente Datenpersistenz durch Twin-Datenbank-Architektur
\end{itemize}

\subsection{Gelerntes}
Während der Entwicklung von SignalReport wurden folgende Kompetenzen erworben:

\begin{itemize}
    \item \textbf{Java 21}: Nutzung moderner Features (Virtual Threads (z.\,B.\ DNS-Benchmark), Pattern Matching, Text Blocks)
    \item \textbf{Web-Entwicklung}: Javalin-Framework, REST-API, Chart.js
    \item \textbf{Datenbank-Programmierung}: H2-Datenbank, SQL, JDBC
    \item \textbf{PDF-Generierung}: OpenPDF und JFreeChart für professionelle Berichte
    \item \textbf{Testing}: JUnit 5, Mockito, Test-Driven Development
    \item \textbf{Software-Architektur}: UML-Diagramme, Design-Patterns, Clean Code
    \item \textbf{Projektmanagement}: Versionskontrolle mit Git, Dokumentation mit LaTeX
\end{itemize}

\subsection{Herausforderungen und Lösungen}
\begin{itemize}
    \item \textbf{Challenge}: IP-Tracking bei dynamischen IPs \\
    \textbf{Lösung}: Separate Tabelle \texttt{ip\_changes} mit Zeitstempeln und Change-Type

    \item \textbf{Challenge}: Maintenance-Fenster mit Über-Mitternacht-Unterstützung \\
    \textbf{Lösung}: Logik zur Erkennung von Zeitfenstern, die über 00:00 Uhr gehen

    \item \textbf{Challenge}: PDF-Generierung mit mehreren Charts und Tabellen \\
    \textbf{Lösung}: OpenPDF statt PDFBox für bessere Schriftarten-Unterstützung

    \item \textbf{Challenge}: Setup-Wizard ohne Kommandozeilen-Abhängigkeit \\
    \textbf{Lösung}: Web-basierte Setup-Seite mit Middleware-Integration

    \item \textbf{Challenge}: Passwort-Sicherheit ohne HTTPS \\
    \textbf{Lösung}: Challenge-Response-Authentifizierung mit SHA-256 und Einmal-Nonces -- Passwort wird nie im Klartext übertragen

    \item \textbf{Challenge (nach Diplomprüfung)}: Datenbank-Korruption durch Windows-Update-Neustart mitten im Schreibvorgang -- Recovery-Tool fand keinen einzigen konsistenten Chunk \\
    \textbf{Lösung}: Twin-Datenbank-Architektur mit synchroner Spiegelung und automatischer Rekonstruktion aus der intakten Datei beim Neustart (siehe Abschnitt \ref{subsec:twindb})
\end{itemize}

\subsection{Mögliche Erweiterungen}
\begin{itemize}
    \item \textbf{HTTPS-Unterstützung}: Verschlüsselte Kommunikation über TLS-Zertifikate
    \item \textbf{Automatischer täglicher PDF-Export}: Cron-Job für regelmäßige Berichte
    \item \textbf{Prometheus Exporter}: Integration in bestehende Monitoring-Tools (Grafana)
    \item \textbf{Multi-Host-Vergleich}: Vergleich von Messungen über mehrere Instanzen
    \item \textbf{Mobile-App}: Native iOS/Android-App für unterwegs
\end{itemize}

\subsection{Fazit}
SignalReport demonstriert, dass eine professionelle Monitoring-Lösung nicht teuer oder komplex sein muss. Mit modernen Java-Technologien und einer durchdachten Architektur entstand eine Anwendung, die:

\begin{itemize}
    \item \textbf{Praktisch nutzbar} ist für echte Provider-Beschwerden
    \item \textbf{Technisch sauber} implementiert ist mit Tests und Dokumentation
    \item \textbf{Erweiterbar} bleibt für zukünftige Anforderungen
    \item \textbf{Lehrreich} war für die persönliche Weiterentwicklung
\end{itemize}

\subsection{Weiterentwicklung nach der Diplomprüfung}
Mit der erfolgreichen Präsentation im April 2026 endete der Diplomlehrgang -- nicht aber das Projekt. SignalReport wird seither als persönliches Open-Source-Vorhaben aktiv weiterentwickelt. Bereits umgesetzte Verbesserungen über den Diplomstand hinaus:

\begin{itemize}
    \item \textbf{Twin-Datenbank-Architektur}: Synchrone Spiegelung und Auto-Recovery (motiviert durch einen realen Korruptionsfall im Produktivbetrieb)
    \item \textbf{H2-Upgrade auf 2.4.240}: Aktuellere DB-Engine mit verbessertem Error-Handling
    \item \textbf{Aufgeräumter Build}: Konsolidierte Maven-Shade-Konfiguration mit korrektem License/Notice-Handling
\end{itemize}

Wer den exakten Stand zum Zeitpunkt der Diplomprüfung begutachten möchte, findet diesen unter dem Git-Tag \texttt{v1.0-diploma} im GitHub-Repository unter \href{https://github.com/matthili/SignalReport/releases/tag/v1.0-diploma}{https://github.com/matthili/SignalReport/releases/tag/v1.0-diploma}.

